\section{Generic-point transfer and concrete algebraic models}
\label{sec:generic-ein-applications}

Theorem~\ref{thm:pure-exp-base} does more than certify isolated special
values.  It provides a concrete transcendence basis on which any finite
algebraic construction may be evaluated without acquiring accidental
relations.  We record this elementary transfer principle before applying it
to Hankel determinants, formal orthogonal polynomials, Grassmannians,
symmetric spectra, and representable matroids.

Throughout this section put $K=\Eexp$.

\begin{proposition}[Generic-point transfer]
\label{prop:generic-point-transfer}
Let $\lambda_1,\ldots,\lambda_s\in\Qbar^\times$ be distinct and put
$y_i=\Ein(\lambda_i)$.  For polynomials
$F_1,\ldots,F_r\in K[X_1,\ldots,X_s]$, the two homomorphisms
\[
 K[T_1,\ldots,T_r]\longrightarrow\C,
 \quad T_j\longmapsto F_j(y_1,\ldots,y_s),
\]
and
\[
 K[T_1,\ldots,T_r]\longrightarrow K[X_1,\ldots,X_s],
 \quad T_j\longmapsto F_j(X_1,\ldots,X_s),
\]
have the same kernel.  The analogous assertion holds for rational maps
after localizing at denominators which do not vanish identically.
\end{proposition}

\begin{proof}
Theorem~\ref{thm:pure-exp-base} says precisely that the evaluation map
$K[X_1,\ldots,X_s]\to\C$, $X_i\mapsto y_i$, is injective.  Composing the
second displayed homomorphism with this injection proves the assertion.
Localization gives the rational version.
\end{proof}

Thus the point $y=(y_1,\ldots,y_s)$ is a generic point of affine
$s$-space over $K$.  The proposition is classical generic-point algebra;
the arithmetic content is the explicit realization of the coordinates as
finite-point $E$-values.

\subsection{A Hankel--Favard transcendence basis}

Put
\begin{equation}\label{eq:ein-moments}
 \mu_m=\Ein(m+1),\qquad
 \Delta_n=\det(\mu_{i+j})_{0\le i,j\le n},
 \qquad \Delta_{-1}=1.
\end{equation}

\begin{theorem}[Hankel determinants and formal Jacobi parameters]
\label{thm:ein-hankel-jacobi}
The countable family
\[
 \Delta_0,\Delta_1,\Delta_2,\ldots
\]
is algebraically independent over $K$, in the finite-subset sense.

Moreover the functional $\mathcal L\in\C[x]^*$ defined by
$\mathcal L(x^m)=\mu_m$ is quasi-definite.  Let its monic formal orthogonal
polynomials be normalized by
\begin{equation}\label{eq:ein-favard-recurrence}
 p_{n+1}(x)=(x-a_n)p_n(x)-b_np_{n-1}(x),
 \qquad p_{-1}=0,\quad p_0=1.
\end{equation}
Then
\[
 \mu_0,a_0,a_1,\ldots,b_1,b_2,\ldots
\]
is algebraically independent over $K$.
\end{theorem}

\begin{proof}
Expansion of $\Delta_n$ at its lower-right entry gives
\begin{equation}\label{eq:hankel-triangular-expansion}
 \Delta_n=\Delta_{n-1}\mu_{2n}
 +R_n(\mu_0,\ldots,\mu_{2n-1})
\end{equation}
for a polynomial $R_n$.  The coefficient $\Delta_{n-1}$ is a nonzero
polynomial in the preceding independent moments.  Hence $\Delta_n$ is
transcendental over $K(\Delta_0,\ldots,\Delta_{n-1})$, and induction proves
the first assertion.  In particular every $\Delta_n$ is nonzero, which is
the quasi-definiteness criterion.

Finite Gram--Schmidt elimination expresses
\[
 (\mu_0,\ldots,\mu_{2N+1})
 \longmapsto
 (\mu_0,a_0,\ldots,a_N,b_1,\ldots,b_N)
\]
rationally.  Conversely, the matrix identity
$\mu_m=\mu_0(J^m)_{00}$ for the finite Jacobi truncation, equivalently the
weighted Motzkin-path expansion, expresses the moments rationally (indeed
polynomially after adjoining $\mu_0$) in the displayed Favard parameters.
The two lists both contain $2N+2$ entries and are birational.  Since the
moment list is algebraically independent, so is the parameter list.  Letting
$N$ vary proves the countable statement.
\end{proof}

The system in Theorem~\ref{thm:ein-hankel-jacobi} is formal rather than
positive-definite: for example $\Delta_1<0$.  It therefore does not define
a self-adjoint Jacobi operator.  Also the Hankel and Favard families must
not be joined into one independent family, since
\begin{equation}\label{eq:hankel-favard-dependence}
 b_n\Delta_{n-1}^2=\Delta_n\Delta_{n-2}.
\end{equation}

\subsection{Pl\"ucker coordinates, spectra, and matroids}

\begin{corollary}[Explicit generic Pl\"ucker data]
\label{cor:ein-pluecker}
Let $1\le r<n$, choose pairwise distinct nonzero algebraic numbers
$\lambda_{ij}$, and form the $r\times n$ matrix
$M=(\Ein(\lambda_{ij}))$.  If $p_I$ denotes its maximal minors, then the
complete relation ideal of the $p_I$ over $K$ is the classical Pl\"ucker
ideal.  Consequently
\[
 \trdeg_K K(p_I)=r(n-r)+1.
\]
The corresponding projective Pl\"ucker ratios have transcendence degree
$r(n-r)$.
\end{corollary}

\begin{proof}
Apply Proposition~\ref{prop:generic-point-transfer} to the maximal-minor
map on a generic $r\times n$ matrix.  Its kernel is the Pl\"ucker ideal and
its image is the affine cone over $\operatorname{Gr}(r,n)$; see, for
example, \cite{BrunsConcaVarbaro2013}.  The dimension statements are the
dimensions of that cone and its projectivization.
\end{proof}

\begin{theorem}[A simple real algebraically independent spectrum]
\label{thm:ein-symmetric-spectrum}
For $1\le i\le j\le n$, choose pairwise distinct positive algebraic
$\lambda_{ij}$ and define the real symmetric matrix
\[
 A_{ij}=A_{ji}=\Ein(\lambda_{ij}).
\]
Then $A$ has $n$ distinct real eigenvalues
$\rho_1,\ldots,\rho_n$, and these eigenvalues are algebraically independent
over $K$.
\end{theorem}

\begin{proof}
The characteristic coefficients of a generic symmetric matrix are
algebraically independent: restriction to diagonal matrices gives the
elementary symmetric polynomials in $n$ independent variables.  The
discriminant is likewise a nonzero polynomial, as diagonal matrices with
distinct entries show.  Proposition~\ref{prop:generic-point-transfer}
therefore makes the characteristic coefficients of $A$ independent and
its discriminant nonzero.  Symmetry makes all roots real, while the
nonzero discriminant makes them distinct.  The splitting field is finite
over the field generated by the $n$ characteristic coefficients, so
\[
 \trdeg_K K(\rho_1,\ldots,\rho_n)=n.
\]
As there are $n$ eigenvalues, they are jointly algebraically independent.
\end{proof}

\begin{corollary}[Representable matroids as $\Ein$-period matroids]
\label{cor:ein-representable-matroid}
Every finite matroid representable over $\Qbar$ occurs as the algebraic
matroid over $K$ of explicit algebraic linear combinations of
$\Ein$-values.  More precisely, if a rank-$m$ representation has columns
$b_e=(b_{1e},\ldots,b_{me})^t$ and
\[
 W_e=\sum_{i=1}^m b_{ie}\Ein(\lambda_i)
\]
for distinct $\lambda_i\in\Qbar^\times$, then for every subset $S$ of the
ground set
\[
 \trdeg_K K(W_e:e\in S)=\operatorname{rank}(b_e:e\in S).
\]
The complete relation ideal is generated by the linear forms
$\sum_e c_eT_e$ with $\sum_e c_eb_e=0$.
\end{corollary}

\begin{proof}
This is Lemma~\ref{lem:linear-image-ein} applied to the representing
matrix and to each column subfamily.
\end{proof}

These corollaries do not introduce new Grassmannian, spectral, or matroid
machinery: each ambient construction is classical.  They give exact
arithmetic realizations by special values and solve no named external open
problem.  The historical priority of these particular $\Ein$ realizations
should be checked separately before publication.
